\chapter{Project overview}

\section{Context and problem statement}
Customer-support teams face growing volumes of repetitive requests, expectations of instant answers and a multiplication of channels (website, voice, e-mail, messaging). Classic chatbots handle open questions poorly, do not always rely on trustworthy sources, and cannot act safely on business systems. A fully human service, on the other hand, is costly and slow for simple requests.

The project targets a platform that:
\begin{itemize}
  \item conducts \textbf{complete, context-aware conversations}, first by text then by voice;
  \item answers from \textbf{approved knowledge} and acknowledges uncertainty;
  \item executes only \textbf{authorized, audited and controlled actions}, enforced by application code;
  \item guarantees a \textbf{human escape hatch}: the customer can always reach a person;
  \item strictly isolates the data of each organization.
\end{itemize}

\section{Vision and product principles}
\begin{encadre}{Vision}
Build an AI-native customer support platform that can conduct complete, context-aware customer conversations over text and voice, answer using trusted business knowledge, perform explicitly authorized actions, and transfer to human agents when appropriate.
\end{encadre}

Eight principles shape the requirements:
\begin{enumerate}
  \item \textbf{Conversation-first}: a conversation is the durable interaction record; channel sessions are delivery mechanisms.
  \item \textbf{Modality-independent intelligence}: text and voice share the same orchestration and authorization model.
  \item \textbf{Grounded answers}: use approved knowledge and business data, and be transparent about uncertainty.
  \item \textbf{Safe action execution}: models \emph{propose}; application services \emph{authorize and execute}.
  \item \textbf{Human escape hatch}: customers can request a human, and policy can escalate automatically.
  \item \textbf{Tenant isolation by default}: data, search, tools, events and analytics.
  \item \textbf{Observable and testable}: every important AI decision, tool execution, handoff and failure is measurable.
  \item \textbf{Provider portability}: model, speech, channel and integration providers are accessed through adapters.
\end{enumerate}

\section{Objectives}
\subsection{Product outcomes}
\begin{itemize}
  \item Reduce time to a correct answer.
  \item Resolve repetitive requests without removing access to human support.
  \item Preserve context across text, voice and channel changes where identity can be safely linked.
  \item Make AI actions auditable and permission-controlled.
  \item Let organizations configure knowledge, tools, policies and escalation behaviour.
  \item Measure \emph{resolution quality}, not only response speed.
\end{itemize}

\subsection{First milestone objective (MVP)}
A website assistant can, end to end:
\begin{enumerate}
  \item accept a customer text message;
  \item maintain a conversation;
  \item retrieve relevant organization knowledge;
  \item invoke one safe, read-only business tool;
  \item answer with source-grounded information;
  \item escalate to a human while preserving context;
  \item emit audit and observability events.
\end{enumerate}
Voice \emph{reuses} the same conversation engine and tool authorization model (release R2).

\subsection{Success measures}
\begin{table}[htbp]
\centering\caption{Product success measures}\label{tab:kpi}
\begin{tabular}{@{}p{4.8cm}p{7.2cm}p{3.6cm}@{}}
\toprule \textbf{Measure} & \textbf{Definition in principle} & \textbf{Target} \\ \midrule
Correct answer rate & Share of answers judged correct on a curated evaluation set. & TBD \\
Groundedness rate & Share of claims supported by a source. & TBD \\
Task completion rate & Share of requests carried to completion (including via a tool). & TBD \\
Human escalation rate & Handoffs relative to volume, segmented by reason. & TBD \\
Median and p95 latency & Response time, for text then for voice. & TBD \\
Tool success rate & Successful calls over authorized calls. & TBD \\
Customer satisfaction & Customer feedback collected after a conversation. & TBD \\
Cost per resolved conversation & Provider and operating cost. & TBD \\
Availability and delivery & Service availability and message-delivery success. & TBD \\
\bottomrule
\end{tabular}
\end{table}
Targets are to be set with stakeholders after a first baseline measurement.

\section{Scope}
\subsection{Included in the MVP (R1)}
\begin{itemize}
  \item Organization, user and access foundation.
  \item Website text chat.
  \item Conversation lifecycle and message persistence.
  \item Basic customer identity.
  \item Knowledge source ingestion and retrieval (RAG).
  \item AI response orchestration.
  \item One read-only business tool (order lookup).
  \item Human inbox and handoff.
  \item Basic ticket creation.
  \item Audit, logs, metrics, traces and usage tracking.
\end{itemize}

\subsection{Later phases}
\begin{itemize}
  \item \textbf{R2}: real-time browser voice, STT/TTS or speech-to-speech depending on the provider, transcript synchronization.
  \item \textbf{R3}: e-mail and external messaging channels, workflow builder, advanced routing and SLA, more connectors, write-capable tools, advanced analytics, evaluation, billing and enterprise governance.
\end{itemize}

\subsection{Explicit non-goals for the first milestone}
\begin{itemize}
  \item Replacing all existing CRM and helpdesk systems.
  \item Arbitrary autonomous changes to customer accounts.
  \item Unreviewed financial transactions.
  \item Supporting every channel and provider at launch.
  \item Training a foundation model.
  \item Splitting every module into a microservice.
\end{itemize}

\section{Stakeholders}
\begin{table}[htbp]
\centering\caption{Stakeholders}\label{tab:stakeholders}
\begin{tabular}{@{}p{4.1cm}p{11.9cm}@{}}
\toprule \textbf{Stakeholder} & \textbf{Main expectations} \\ \midrule
End customers & Fast and correct answers, access to a human, protected data. \\
Support agents & Conversation takeover with context, a clear work queue, efficient tools. \\
Support leads & Control of queues, policies, quality and SLAs. \\
Organization administrators & Configuration of channels, knowledge, AI agents, integrations and access. \\
Platform operators & Reliability, cost, abuse control and system health. \\
Security and privacy & Tenant isolation, data minimization, auditability. \\
Business-system owners & Controlled, limited access to their APIs. \\
\bottomrule
\end{tabular}
\end{table}

\section{General assumptions and constraints}
\subsection{Assumptions}
\begin{itemize}
  \item A small team initially builds a \textbf{modular monolith}; services are extracted only when scaling, fault isolation or deployment cadence justify it.
  \item Product design, infrastructure, security and automated tests are part of delivery, not later phases.
  \item Story points are team-specific estimates, not time conversions.
  \item Sprints are assumed to last two weeks; the sequence is provisional.
\end{itemize}

\subsection{Constraints}
\begin{itemize}
  \item AI model output is \textbf{untrusted input}: authorization and business invariants are enforced by application code.
  \item The core domain depends on no specific model or messaging provider; provider details live in adapters.
  \item No raw audio or sensitive customer data is stored without a clear product need, consent or legal basis, access policy and retention policy.
  \item Every request or job carries a trace identifier and tenant context.
\end{itemize}
